\section{Reduced costs}
\begin{defn}[Potentials and reduced costs \cite{ahuja1993network}]\label{def:reduced_cost}
For each node $i \in V$ we associate a real number $\pi(i)$ as the \textbf{potential} of node $i$. For a given set of node potentials $\pi$, we define the \textbf{reduced cost} of arc $(i,j) \in E$ as $c^\pi_{ij} = c_{ij} - \pi(i) + \pi(j) $ 
\end{defn}

We show some useful correlations between networks with cost function $c$ and networks with nodes associated with node potentials $\pi$ and arc cost function $c^\pi$. Let $z(\pi) = \sum_{(i,j) \in E} c^\pi_{ij}x_{ij}$ denote an objective function for a minimum cost flow.

\begin{lemma}[Property 2.4 in \cite{ahuja1993network}]
\label{reduced_equiv}
Minimum cost flow problems with arc costs $c_{ij}$ and $c_{ij}^\pi$ have the same optimal solutions. Moreover, $z(\pi) = z(0) - \pi b$.
\end{lemma}
\begin{proof}[Proof (see{ \cite[p.~43]{ahuja1993network}})]
Consider an initial node potential $\pi_s$. Let us define a new potential $\pi_s'$ by increasing the potential at a specific node $k$ by $\alpha$. That is, $\pi_s'(k) = \pi_s(k) + \alpha$, and $\pi_s'(i) = \pi_s(i)$ for all $i \neq k$. 

Recall the flow conservation constraint for any node $i$: $\sum_j x_{ij} - \sum_j x_{ji} = b(i)$. Thus evaluating the objective function under the new potential yields:
\begin{align*}
    z(\pi_s') = \sum_{(i,j)} c^{\pi_s'}_{ij} x_{ij} 
    &= \sum_{(i,j)} c^{\pi_s}_{ij} x_{ij} - \sum_i \alpha x_{ki} + \sum_i \alpha x_{ik} \\
    &= z(\pi_s) - \alpha b(k)
\end{align*}

Since $\alpha b(k)$ is a constant with respect to the flow variables $x_{ij}$, the optimal solution remains unchanged. 

Therefore, if we start with a zero potential vector ($0$) and incrementally update it to reach an arbitrary potential vector $\pi$, we preserve optimality and obtain the equality $z(\pi) =z(0)- \sum_k \pi(k)\cdot b(k)=  z(0) - \pi b$.
\end{proof}

From the definition of reduced costs, the following lemma follows directly:
\begin{lemma}[Property 9.2 in \cite{ahuja1993network}]
\label{reducedproperties2}
The following properties occur
\begin{enumerate}
    \item For any directed path $P$ from node $k$ to node $l$, $\sum_{(i,j)\in P} c_{ij}^\pi = \sum_{(i,j)\in P} c_{ij} - \pi(k) + \pi(l)$
    \item For any directed cycle $W$, $\sum_{(i,j)\in W} c_{ij}^\pi = \sum_{(i,j)\in W} c_{ij}$
\end{enumerate}
\end{lemma}

\section{Optimality conditions}\label{sec:optimality_conditions_flow}
Now we introduce the optimality conditions for residual networks, which will be used to prove the complementary slackness optimality conditions.

To prove optimality conditions theorem, let us first introduce the following theorem, which we state without proof.
\begin{theorem}[Augmenting cycle theorem \cite{ahuja1993network}]
\label{thm:augmenting_cycle}
Let $x$ and $x^*$ be any two feasible solutions of a network flow problem. Then $x$ equals $x^*$ plus the flow on at most $m$ directed cycles in $G(x^*)$. Furthermore, the cost of $x$ equals the cost of $x^*$ plus the cost of flow on these augmenting cycles.
\end{theorem}


\begin{theorem}[Negative cycle optimality conditions \cite{ahuja1993network}]
\label{optimality1}
A feasible solution $x^*$ is an optimal solution of the minimum cost flow problem if and only if it satisfies the negative cycle optimality conditions: namely, the residual network $G(x^*)$ contains no negative (directed) cycle.
\end{theorem}

\begin{proof}[Proof (see {\cite[p.~307]{ahuja1993network}})]

($\Rightarrow$) Suppose that there $G(x)$ contains a negative-cost cycle. By pushing a positive flow along this cycle, we obtain a new feasible flow $x'$ with a strictly lower cost. Therefore, $x$ cannot be optimal. 

($\Leftarrow$) Suppose that $x$ is a feasible flow and the residual network $G(x)$ contains no negative-cost cycles. Let $x^*$ be an optimal flow such that $x^* \neq x$. By \Cref{thm:augmenting_cycle} the flow difference $x^*-x$ can be decomposed to at most $m$ augmenting cycles in $G(x)$. Due to the inital assumption, none of those cycles are negative-cost. Therefore, we have $c \cdot (x^* - x) \geq 0$, or equvalently $cx^* \geq cx$. Since $x^*$ is an optimal have $cx^* \leq cx$. Thus $cx^*= cx$, which implies that $x$ is also an optimal flow. 

\end{proof}


\begin{theorem}[Reduced cost optimality conditions \cite{ahuja1993network}]
\label{red_cost_opt_cond}
A feasible solution $x^*$ is an optimal solution of the minimum cost flow problem if and only if some set of node potentials $\pi$ satisfy the following reduced cost optimality conditions:
\[
c_{ij}^\pi \geq 0 \quad \text{for every arc } (i,j) \text{ in } G(x^*).
\]
\end{theorem}

\begin{proof}[Proof (see {\cite[p.309]{ahuja1993network}})]

($\Leftarrow$)
Suppose that a feasible flow $x^*$ satisfies $c_{ij}^\pi \geq 0$ for every arc $(i,j)$ in $G(x^*)$ for some $\pi$. Then, for each directed cycle $W$ in $G(x^*)$, we have $\sum_{(i,j)\in W} c_{ij}^\pi \geq 0$. By \Cref{reducedproperties2}, we obtain $\sum_{(i,j)\in W} c_{ij}^\pi = \sum_{(i,j)\in W} c_{ij} \geq 0$, meaning $G(x^*)$ has no negative cycles. Therefore, by \Cref{optimality1}, $x^*$ is optimal.

($\Rightarrow$) 
Assume that $x^*$ is an optimal solution. By \Cref{optimality1}, $G(x^*)$ contains no negative cycles. Therefore, we can compute the shortest path distances from an arbitrary node $s$ to all other nodes in $G(x^*)$, using $c_{ij}$ as the weight for each arc $(i,j)$. Let $d(i)$ denote the distance from node $s$ to node $i$. By the triangle inequality, we have $d(j) \leq d(i) + c_{ij}$ for all $(i,j)$ in $G(x^*)$. Equivalently, this can be rewritten as $c_{ij} - (-d(i)) + (-d(j)) \geq 0$, which yields $c_{ij}^\pi \geq 0$ if we define $\pi = -d$.
\end{proof}

\subsection{Complementary slackness optimality conditions}
We previously established the conditions for residual networks. Now, we will use those conditions to get the conditions for the original networks.

\begin{theorem}[Complementary slackness optimality conditions \cite{ahuja1993network}]
\label{comp_slack}
    A feasible solution $x^*$ is an optimal solution of minimal cost flow problem if and only if for some set of node potentials $\pi$, the reduced costs and flow values satisfy the following complementary slackness optimality conditions for every arc $(i,j) \in E$:
\begin{align*}
    c_{ij}^\pi > 0 &\Rightarrow x_{ij}^* = 0 \\
    0 < x_{ij}^* < u_{ij} &\Rightarrow c_{ij}^\pi = 0 \\
    c_{ij}^\pi < 0 &\Rightarrow x_{ij}^* = u_{ij}
\end{align*}
\end{theorem}

\begin{proof}[Proof (see {\cite[p.310]{ahuja1993network}})]
We show a proof of equivalence to the \Cref{red_cost_opt_cond} in one direction for simplicity. 
\begin{itemize}
    \item Suppose that $c_{ij}^\pi > 0$ and $x^*_{ij} > 0$. Then, the residual network contains $(j,i)$. That contradicts
    \Cref{red_cost_opt_cond}, since $c_{ji}^\pi = -c_{ij}^\pi < 0$.

    \item Suppose that $0 < x_{ij}^* < u_{ij}$. The residual networks contains both $(i,j)$ and $(j,i)$. By \Cref{red_cost_opt_cond}, ${c_{ij}^\pi \geq 0}$ and ${c_{ji}^\pi \geq 0}$, thus $c_{ij}^\pi = c_{ji}^\pi =0$. 

    \item Suppose that $c_{ij}^\pi < 0$ and $x^*_{ij} < u_{ij}$. Then the residual network contains $(i,j)$. That contradicts \Cref{red_cost_opt_cond}, since $c_{ij}^\pi < 0$.

    
\end{itemize}

    The converse is left to the reader as an exercise.
\end{proof}